% ======================================================================
\chapter{Materiais e métodos}
\label{cap:metodologia}
% ======================================================================
\instrucao{Descreva a montagem com um desenho, liste os instrumentos com a resolução de cada um e conte o procedimento no passado, com detalhe suficiente para outra pessoa repetir o experimento.}

\section{Montagem experimental}

A \autoref{fig:montagem} mostra a montagem. Um tubo de acrílico de \qty{1.0}{\metre} contém água, cujo nível é controlado por um reservatório móvel ligado ao fundo do tubo. Um alto-falante, alimentado por um gerador de sinais, fica logo acima da boca do tubo.

\begin{figure}[htb]
  \caption{Esquema do tubo de ressonância com o nível da água ajustável.}
  \label{fig:montagem}
  \centering
  \begin{tikzpicture}[>=Stealth, scale=0.9]
    % tubo
    \draw[thick] (0,0) -- (0,6) (1.2,0) -- (1.2,6);
    \draw[thick] (0,0) -- (1.2,0);
    % água
    \fill[blue!15] (0.02,0.02) rectangle (1.18,2.4);
    \draw[blue!60!black, decorate, decoration={snake, amplitude=0.6pt, segment length=6pt}]
      (0.02,2.4) -- (1.18,2.4);
    % envoltória da onda estacionária na coluna de ar
    \draw[gray, domain=2.4:6.3, samples=60, smooth]
      plot ({0.6 + 0.45*sin(deg((\x-2.4)*pi/2.6))}, \x);
    \draw[gray, domain=2.4:6.3, samples=60, smooth]
      plot ({0.6 - 0.45*sin(deg((\x-2.4)*pi/2.6))}, \x);
    % alto-falante
    \fill[gray!60] (0.35,6.6) rectangle (0.85,6.85);
    \draw (0.2,6.9) -- (0.35,6.6) -- (0.85,6.6) -- (1.0,6.9) -- cycle;
    \draw (0.6,6.9) -- (0.6,7.4) -- (3.4,7.4) node[right, align=left] {gerador\\de sinais};
    % reservatório
    \draw[thick] (3,1.2) rectangle (4,3.4);
    \fill[blue!15] (3.02,1.22) rectangle (3.98,2.4);
    \draw[thick] (0.6,0) -- (0.6,-0.5) -- (3.5,-0.5) -- (3.5,1.2);
    \draw[<->] (4.3,1.6) -- (4.3,3.0) node[midway, right] {sobe e desce};
    % cota L
    \draw[<->] (-0.4,6) -- (-0.4,2.4) node[midway, left] {$L_n$};
  \end{tikzpicture}

  \legend{Fonte: elaborada pelo autor.}
\end{figure}

\section{Materiais}

\begin{alineas}
  \item tubo de acrílico com diâmetro interno de \qty{38}{\milli\metre};
  \item gerador de sinais senoidais com resolução de \qty{1}{\hertz};
  \item alto-falante de \qty{8}{\ohm};
  \item trena com resolução de \qty{1}{\milli\metre};
  \item termômetro com resolução de \qty{0.5}{\celsius}.
\end{alineas}

\section{Procedimento}

Com o gerador em \qty{\frequencia}{\hertz}, o nível da água foi abaixado lentamente a partir da boca do tubo. Cada ressonância foi identificada pelo máximo de intensidade sonora, e a posição do nível foi lida na trena. Cada posição foi medida cinco vezes, com aproximação por cima e por baixo, e a média foi adotada como $L_n$. A temperatura do ar foi registrada no início e no fim da série.
